% Part: first-order-logic
% Chapter: model-theory
% Section: partial-iso

\documentclass[../../../include/open-logic-section]{subfiles}

\begin{document}

\olfileid[hi]{mod}{bas}{dlo}
\section{सघन रैखिक क्रम}

\begin{defn}
  कोई \emph{अंतबिंदु-रहित सघन रैखिक क्रम} ऐसी !!a{structure}~$\Struct{M}$
  है जो केवल एक द्वि-स्थानीय !!{predicate}~$<$ वाली !!{language} के लिए है
  और निम्न वाक्यों को संतुष्ट करती है:
  \begin{enumerate}
  \item $\lforall[x][\lnot x < x]$;
  \item $\lforall[x][\lforall[y][\lforall[z][(x < y \lif (y < z \lif x
    <z ))]]]$;
  \item $\lforall[x][\lforall[y][(x< y \lor \eq[x][y] \lor y < x)]]$;
  \item $\lforall[x][\lexists[y][x < y]]$;
  \item $\lforall[x][\lexists[y][y < x]]$;
  \item $\lforall[x][\lforall[y][(x < y \lif \lexists[z][(x < z \land
        z < y))]]]$.
 \end{enumerate}
\end{defn}

\begin{thm}\ollabel{thm:cantorQ}
  कोई भी दो !!{enumerable} अंतबिंदु-रहित सघन रैखिक क्रम समरूपी होते हैं।
\end{thm}

\begin{proof}
  मान लें कि $\Struct{M_1}$ और $\Struct{M_2}$ ऐसे !!{enumerable}
  अंतबिंदु-रहित सघन रैखिक क्रम हैं जिनमें ${<_1} = \Assign{<}{M_1}$ और ${<_2} =
  \Assign{<}{M_2}$, तथा $\PIso{I}$ उनके बीच सभी आंशिक समरूपताओं का समुच्चय
  है। $\PIso{I}$ अरिक्त है, क्योंकि कम-से-कम $\emptyset \in \PIso{I}$। हम
  दिखाते हैं कि $\PIso{I}$ आगे-पीछे गुणधर्म को संतुष्ट करता
  है। तब $\Struct{M_1} \iso[p] \Struct{M_2}$, और प्रमेय
  \olref[pis]{thm:p-isom1} से मिलता है।

  यह दिखाने के लिए कि $\PIso{I}$ आगे गुणधर्म को संतुष्ट करता है, कोई $p \in
  \PIso{I}$ लें और $p(a_i) = b_i$ मानें, जहाँ $i = 1$, \dots,~$n$।
  व्यापकता की हानि के बिना मान लें कि $a_1 <_1 a_2 <_1 \cdots <_1
  a_n$।
  कोई $a \in \Domain{M_1}$ दिया हो, तो $b \in \Domain{M_2}$ को इस प्रकार
  खोजें:
  \begin{enumerate}
  \item यदि $a <_1 a_1$, तो $b \in \Domain{M_2}$ ऐसा लें कि $b <_2
    b_1$;
  \item यदि $a_n <_1 a$, तो $b \in \Domain{M_2}$ ऐसा लें कि $b_n <_2 b$;
 \item यदि $a_i <_1 a <_1 a_{i+1}$ किसी $i$ के लिए, तो $b \in
   \Domain{M_2}$ ऐसा लें कि $b_i <_2 b <_2 b_{i+1}$।
  \end{enumerate}
  वांछित गुणधर्म वाला $b$ खोजना हमेशा संभव है, क्योंकि $\Struct{M_2}$ कोई
  अंतबिंदु-रहित सघन रैखिक क्रम है। $q = p \cup \{ \langle a, b \rangle \}$
  परिभाषित करें; तब $q \in \PIso{I}$, $p$ का वांछित विस्तार है। इससे
  आगे गुणधर्म स्थापित होता है। पीछे गुणधर्म का प्रमाण समान है। अतः
  $\Struct{M_1} \iso[p]
  \Struct{M_2}$; \olref[pis]{thm:p-isom1} से $\Struct{M_1} \iso
  \Struct{M_2}$।
\end{proof}

\begin{prob}
  यह सत्यापित करके \olref[mod][bas][dlo]{thm:cantorQ} का प्रमाण पूरा कीजिए
  कि $\PIso{I}$ पीछे गुणधर्म को संतुष्ट करता है।
\end{prob}

\begin{rem}
  मान लें कि $\Struct{S}$ कोई भी !!{enumerable} अंतबिंदु-रहित सघन रैखिक
  क्रम है। तब (\olref{thm:cantorQ} से) $\Struct{S} \iso
  \Struct{Q}$, जहाँ
  $\Struct{Q} = (\Rat, <)$ वह !!{enumerable} सघन रैखिक क्रम है जिसका प्रांत
  परिमेय संख्याओं का समुच्चय $\Rat$ है। अब \olref[thm]{remark:R} की
  !!{structure}~$\Struct{R} =
  (\Real, <)$ पर फिर विचार करें। हमने देखा कि
  कोई !!a{enumerable} !!{structure}~$\Struct{S}$ ऐसी है कि $\Struct{R}
  \elemequiv \Struct{S}$। किंतु $\Struct{S}$ कोई !!a{enumerable}
  अंतबिंदु-रहित सघन रैखिक क्रम है, इसलिए वह !!{structure}~$\Struct{Q}$ के
  समरूपी (और अतः प्राथमिकतः तुल्य) है। प्राथमिक तुल्यता की संक्रामकता से
  $\Struct{R} \elemequiv
  \Struct{Q}$। (इसे हम उसी आगे-पीछे तर्क से सीधे
  $\Struct{R} \iso[p] \Struct{Q}$ स्थापित करके भी दिखा सकते थे।)
\end{rem}
\end{document}
